% !TeX root=../main.tex
\chapter{روش طراحی و پیاده‌سازی}
\label{chap:method}

\section{مقدمه}

هدف این پروژه ساخت ابزاری است که فضای حالت زمان‌دار صادرشده از افرا را دریافت کند، تعامل‌های خارج از مجموعهٔ مورد توجه کاربر را پنهان سازد، حالت‌ها را بر اساس شبیه‌سازی دوسویهٔ ضعیف زمان‌دار افراز کند و سامانهٔ خارج‌قسمتی حاصل را تحویل دهد. ورودی الگوریتم یک سامانهٔ انتقال متناهی است، نه کد منبع ربکا؛ بنابراین، تولید فضای حالت همچنان بر عهدهٔ زنجیرهٔ ابزار ربکا و افراست. پیاده‌سازی حاصل با نام \lr{awtr} در مخزن پروژه نگهداری می‌شود \cite{AfraWeakTimedReduction2026}.

روش انجام کار ترکیبی از طراحی الگوریتم، مهندسی نرم‌افزار و ارزیابی تجربی است. ابتدا قرارداد واقعی فایل خروجی افرا از روی تولیدکننده و خواننده‌های رسمی آن استخراج شد. سپس ساختار الگوریتم بر پایهٔ پالایش افراز—همان ایده‌ای که گروته و واندراخر برای شبیه‌سازی دوسویهٔ انشعابی به کار بردند \cite{GrooteVaandrager1990}—برای حالت زمان‌دار تطبیق داده شد، و نحوهٔ برخورد با تأخیرها مطابق ویژگی‌های تعریف \ref{def:tts-properties} تعیین شد. پس از ساخت افراز، \gls{QuotientSystem} تولید و با وارسی‌گری مستقل بررسی شد. در پایان، رفتار ابزار با آزمون‌های واحد، \glspl{EndToEndTesting} و ورودی‌های ارزیابی سنجیده شد. این فصل روش و تصمیم‌های طراحی را شرح می‌دهد و فصل \ref{chap:results} نتایج ثبت‌شده را گزارش می‌کند.

\section{صورت مسئله، ورودی و خروجی}

فرض کنید سامانهٔ انتقال زمان‌دار متناهی پس از خواندن فایل ورودی به‌صورت
\[
T=(S,s_0,\mathit{Act}\cup\{\tau\}\cup\mathbb{N}_{>0},\rightarrow)
\]
نمایش داده شود. $S$ مجموعهٔ حالت‌ها، $s_0$ حالت آغازین، $\mathit{Act}$ مجموعهٔ تعامل‌های نام‌دار و $d\in\mathbb{N}_{>0}$ برچسب گذر دقیقاً $d$ واحد زمان است. کاربر در هر اجرا زیرمجموعه‌ای مانند $O\subseteq\mathit{Act}$ را به‌عنوان تعامل‌های قابل مشاهده تعیین می‌کند. همهٔ تعامل‌های دیگر به $\tau$ نگاشت می‌شوند و سپس افراز پایدار متناظر با تعریف \ref{def:weak-timed-bisimulation} محاسبه می‌شود.

خروجی اصلی فقط یک پاسخ بله یا خیر نیست. ابزار باید سه شیء مرتبط را تولید کند: نگاشت هر حالت ورودی به ردهٔ هم‌ارزی، سامانهٔ انتقال خارج‌قسمتی، و رابطه‌ای که نشان دهد مدل ورودی و خروجی مطابق تعریف انتخاب‌شده رفتار یکسان دارند. علاوه بر این، برای تکرارپذیری هر اجرا، مجموعهٔ تعامل‌های قابل مشاهده، هش ورودی، نسخهٔ ابزار، شناسهٔ کامیت، زمان اجرا و اندازه‌های پیش و پس از کاهش ثبت می‌شوند \cite{AfraWeakTimedReduction2026}.

دامنهٔ روش به زمان گسسته و تأخیرهای صحیح نامنفی محدود است. ورودی‌های زمان چگال، معنای \lr{FTTS}، وارسی شبیه‌سازی دوسویهٔ قوی و کامپایل مستقیم کد ربکا در این پروژه پیاده‌سازی نشده‌اند. همچنین، هدف ساخت مدلی کوچک‌تر و وارسی‌شده است؛ مینیمال یا یکتابودن مدل خروجی ادعا نمی‌شود.

\section{جریان پردازش}

شکل \ref{fig:reduction-pipeline} مسیر داده را از فایل افرا تا مصنوعات خروجی نشان می‌دهد. مرز میان خواندن فایل و هستهٔ الگوریتم صریح نگه داشته شده است؛ ازاین‌رو، تغییر منبع داده، منطق \gls{Hiding}، \gls{PartitionRefinement} یا ساخت خارج‌قسمت را تغییر نمی‌دهد.

\begin{figure}[htbp]
\centering
\begin{tikzpicture}[
stage/.style={draw, rounded corners, align=center, minimum width=3.1cm, minimum height=0.9cm},
flow/.style={->, thick}
]
\node[stage] (input) at (-4.2,0) {خروجی \lr{TTS} افرا\\\lr{.statespace}};
\node[stage] (read) at (0,0) {خواندن و\\اعتبارسنجی};
\node[stage] (hide) at (4.2,0) {حذف حالات دور از دسترس\\و پنهان‌سازی};
\node[stage] (unit) at (4.2,-2.0) {پالایش تأخیرهای\\چندواحدی};
\node[stage] (weak) at (0,-2.0) {ساخت رابطهٔ\\گذار ضعیف};
\node[stage] (partition) at (-4.2,-2.0) {پالایش افراز تا\\رسیدن به پایداری};
\node[stage] (quotient) at (-4.2,-4.0) {ساخت و فشرده‌سازی\\مدل خارج‌قسمتی};
\node[stage] (verify) at (0,-4.0) {وارسی مستقل\\رابطهٔ حاصل};
\node[stage] (output) at (4.2,-4.0) {مدل کاهش‌یافته، افراز\\و سنجه‌های اجرا};
\draw[flow] (input) -- (read);
\draw[flow] (read) -- (hide);
\draw[flow] (hide) -- (unit);
\draw[flow] (unit) -- (weak);
\draw[flow] (weak) -- (partition);
\draw[flow] (partition) -- (quotient);
\draw[flow] (quotient) -- (verify);
\draw[flow] (verify) -- (output);
\end{tikzpicture}
\caption{جریان پردازش در ابزار کاهش فضای حالت زمان‌دار.}
\label{fig:reduction-pipeline}
\end{figure}

این تفکیک دو فایده دارد. نخست، هر مرحله قرارداد مشخصی دارد و می‌توان آن را جداگانه آزمود. دوم، وارسی‌گر خروجی از پیاده‌سازی پالایش افراز استفادهٔ مجدد نمی‌کند؛ در نتیجه، خطای مشترک در محاسبهٔ رابطهٔ ضعیف و وارسی آن به‌سادگی پنهان نمی‌شود.

\section{خواندن خروجی افرا}

\subsection{قرارداد فایل \lr{TTS}}

ساختار فایل \lr{.statespace} با بررسی تولیدکنندهٔ رسمی \lr{RMC}، طرح دادهٔ افرا و دو خوانندهٔ رسمی فضای حالت تعیین شد \cite{AfraGitHub2026,RMCGitHub2026,StateSpaceTransformerGitHub2026,StateSpaceAnalysisGitHub2026}. قواعدی که مستقیماً بر پیاده‌سازی اثر گذاشتند در جدول \ref{tab:afra-contract} آمده‌اند.

\begin{table}[htbp]
\caption{بخش مورد استفاده از قرارداد خروجی فضای حالت افرا}
\label{tab:afra-contract}
\centering
\onehalfspacing
\begin{tabularx}{0.94\textwidth}{|r|X|}
\hline
جزء فایل & تفسیر در ابزار \\
\hline
\lr{state/@id} & شناسهٔ حالت به‌صورت رشته‌ای و بدون تفسیر ساختار داخلی آن نگهداری می‌شود. \\
\hline
نخستین عنصر \lr{state} & حالت آغازین است؛ گزینهٔ \lr{--initial-state} امکان جایگزینی صریح آن را فراهم می‌کند. \\
\hline
\lr{messageserver} & ویژگی \lr{title} هویت پایهٔ تعامل است؛ صورت \lr{owner.title} نیز برای رفع ابهام قابل انتخاب است. \\
\hline
\lr{time/@value} & مدت سپری‌شده است و باید یک عدد صحیح نامنفی باشد. \\
\hline
\lr{executionTime} و \lr{shift} & برای ردیابی نگه داشته می‌شوند، اما در مدت گذار دخالت داده نمی‌شوند. \\
\hline
\end{tabularx}
\end{table}

یک ناسازگاری ظاهری در خروجی واقعی وجود دارد: تولیدکنندهٔ \lr{RMC} عنصر ریشهٔ \lr{transitionsystem} را باز می‌کند، اما برچسب پایانی آن را نمی‌نویسد. بنابراین، فایل واقعی از دید یک تجزیه‌گر سخت‌گیر، \lr{XML} خوش‌ساخت نیست. خوانندهٔ این پروژه انتهای فایل را بررسی می‌کند و فقط در صورت نبودن برچسب پایانی، آن را به جریان ورودی اضافه می‌کند. به این ترتیب فایل بدون ویرایش دستی مصرف می‌شود و فایل‌های دارای برچسب پایانی نیز پذیرفته می‌شوند \cite{RMCGitHub2026,StateSpaceAnalysisGitHub2026,AfraWeakTimedReduction2026}.

تنها خروجی‌های \lr{TTS} برای کاهش زمانی مناسب‌اند؛ زیرا در این حالت، گذر زمان به‌صورت عناصر مستقل \lr{time} ثبت می‌شود. در \lr{FTTS} زمان در گذارهای پیام ادغام می‌شود و یال تأخیر مستقلی برای الگوریتم حاضر وجود ندارد. خواننده همچنین شناسهٔ تکراری حالت، مقصد یا مبدأ تعریف‌نشده، زمان منفی یا اعشاری، پیام فاقد ویژگی‌های لازم، و وجود هم‌زمان برچسب پیام و زمان روی یک گذار را رد می‌کند. پردازش \lr{DTD} و موجودیت خارجی نیز غیرفعال شده است.

\subsection{مرز منبع داده}

خوانندهٔ فایل یکی از پیاده‌سازی‌های رابط \lr{TransitionSystemSource} است. این رابط تنها یک عمل برای ساخت نمایش خام سامانه دارد. نمایش خام، نام فرستنده، مالک و عنوان پیام و نیز اطلاعات زمانی جانبی را حفظ می‌کند. سپس مرحلهٔ جداگانه‌ای آن را به سامانهٔ انتقال مورد استفادهٔ الگوریتم تبدیل می‌کند. یک منبع درون‌حافظه‌ای نیز همین رابط را پیاده‌سازی می‌کند و آزمون قرارداد نشان می‌دهد هر دو منبع به مدل دامنه و نتیجهٔ کاهش یکسان می‌رسند. این مرز امکان افزودن یک آداپتور درون‌فرایندی برای افرا را فراهم می‌کند، بی‌آنکه هستهٔ کاهش به کلاس‌های رابط کاربری یا \lr{XML} وابسته شود \cite{AfraWeakTimedReduction2026}.

\section{انتخاب مشاهده‌ها و پیش‌پردازش}

مجموعهٔ تعامل‌های قابل مشاهده در زمان اجرا و با گزینهٔ \lr{--observable} تعیین می‌شود. یک نام ساده مانند \lr{getSense} با عنوان پیام در هر مالک تطبیق داده می‌شود؛ صورت مقید مانند \lr{controller.getSense} فقط همان مالک را می‌پذیرد. تطبیق به بزرگی و کوچکی حروف حساس است. این تصمیم از ادغام ناخواستهٔ دو شناسهٔ متمایز جلوگیری می‌کند و با هویت ثبت‌شده در خروجی افرا سازگار است \cite{StateSpaceAnalysisGitHub2026,StateSpaceTransformerGitHub2026}.

اگر $\alpha$ یک تعامل نام‌دار و $O$ مجموعهٔ انتخابی کاربر باشد، تابع پنهان‌سازی به‌صورت زیر عمل می‌کند:
\[
h_O(\alpha)=
\begin{cases}
\alpha & \alpha\in O,\\
\tau & \alpha\notin O.
\end{cases}
\]
گذر بدون برچسب و گذر زمان صفر نیز به $\tau$ تبدیل می‌شوند. عنوان‌هایی که افرا با پیشوند \lr{tau=>} به‌عنوان ادامهٔ اجرای یک پیام ثبت کرده است، با یک نام ساده قابل مشاهده نمی‌شوند. پس از پنهان‌سازی، فقط زیرسامانهٔ قابل دسترس از حالت آغازین نگه داشته می‌شود. ابزار نام‌های قابل مشاهده‌ای را که با هیچ گذار ورودی تطبیق ندارند نیز گزارش می‌کند؛ زیرا یک خطای تایپی در این گزینه می‌تواند ناخواسته همهٔ تعامل‌ها را پنهان کند.

\section{چرا یال‌های تأخیر به گام‌های یک‌واحدی شکسته می‌شوند}
\label{sec:unit-splitting}
\label{sec:time-semantics-decision}

مهم‌ترین تصمیم معنایی این پروژه، شکستن هر یال تأخیر با مدت $d$ به $d$ یال یک‌واحدی است. این تصمیم یک بهینه‌سازی یا سلیقهٔ پیاده‌سازی نیست؛ بدون آن، الگوریتم افراز نمی‌تواند تعریف \ref{def:weak-timed-bisimulation} را روی خروجی افرا محاسبه کند. دلیل آن به فاصلهٔ میان دو چیز بازمی‌گردد: سامانهٔ انتقال زمان‌داری که مدل را توصیف می‌کند، و گراف متناهی‌ای که افرا از آن ذخیره می‌کند.

\subsection{مسئله: حالت‌هایی که در گراف ذخیره‌شده شریکی ندارند}

دو سامانهٔ زیر را در نظر بگیرید که در آن‌ها $a$ تنها کنش قابل مشاهده است:
\[
\begin{array}{rcl}
A:\quad q_0 & \xrightarrow{2} & q_1\xrightarrow{3}q_2\xrightarrow{a}q_3,\\[2mm]
B:\quad p_0 & \xrightarrow{5} & p_1\xrightarrow{a}p_2.
\end{array}
\]
هر دو سامانه دقیقاً یک رفتار قابل مشاهده دارند: پس از گذشت پنج واحد زمان، کنش $a$ رخ می‌دهد. طبق تعریف \ref{def:weak-timed-bisimulation} این دو سامانه شبیه‌سازی دوسویهٔ ضعیف زمان‌دار دارند، زیرا رابطهٔ گذار ضعیف بر اساس \emph{وجود یک اجرا} با مدت کل معین تعریف شده است، نه بر اساس وجود یک یال منفرد؛ در $A$ اجرایی از $q_0$ به $q_2$ با مدت کل ۵ و بدون کنش قابل مشاهده وجود دارد، پس $q_0\overset{5}{\Rightarrow}q_2$ برقرار است و $p_0\overset{5}{\Rightarrow}p_1$ نیز.

اکنون فرض کنید بخواهیم چنین رابطه‌ای را فقط میان حالت‌های نوشته‌شده در همین دو گراف بسازیم. حالت $q_1$، یعنی لحظهٔ «دو واحد گذشته است»، در $A$ صریحاً وجود دارد؛ زیرا $q_0\overset{2}{\Rightarrow}q_1$ برقرار است، رابطه باید برای $q_1$ شریکی در $B$ بیابد. اما در $B$ هیچ حالتی متناظر با لحظهٔ «دو واحد گذشته است» نوشته نشده است: تنها گذار تأخیری موجود، یک یال پنج‌واحدی است و $p_0\overset{2}{\Rightarrow}p'$ برای هیچ $p'$ از حالت‌های $\{p_0,p_1,p_2\}$ برقرار نیست. در نتیجه $q_1$ با هیچ حالتی جفت نمی‌شود و ساخت رابطه شکست می‌خورد—درحالی‌که دو سامانه طبق تعریف هم‌ارزند.

\subsection{راه‌حل: استفاده از ویژگی پیوستگی}

این تناقض ظاهری از یک اشتباه در سطح تحلیل می‌آید. مطابق تعریف \ref{def:tts-properties}، رابطهٔ گذار یک سامانهٔ انتقال زمان‌دار ویژگی \emph{پیوستگی} دارد: اگر $p_0\xrightarrow{5}p_1$ و $5=2+3$، آنگاه پیکربندی $p'$ وجود دارد که
\[
p_0\xrightarrow{2}p'\xrightarrow{3}p_1 .
\]
یعنی لحظهٔ میانی «دو واحد گذشته است» در \emph{سامانهٔ انتقال زمان‌دار} $B$ حتماً وجود دارد؛ آنچه وجود ندارد، \emph{نمایش صریح} آن در فایل خروجی افراست. افرا فقط تعدادی یال تأخیر با مدت‌های معین را ثبت می‌کند و پیکربندی‌های میانی را نمی‌نویسد، زیرا برای تحلیل‌های متداول خود به آن‌ها نیاز ندارد.

برای آنکه الگوریتم افراز بتواند روی همین گراف متناهی کار کند، باید آن پیکربندی‌های میانی را که تعریف به آن‌ها نیاز دارد، صریحاً بسازیم. در حالت کلی این کار شدنی نیست: در زمان چگال میان هر دو لحظه بی‌نهایت لحظهٔ دیگر وجود دارد. اما در این پروژه، همان‌طور که در بخش پیش‌تر آمد، دامنهٔ زمان $\mathbb{T}=\mathbb{N}$ است و همهٔ مدت‌های ثبت‌شده اعداد صحیح‌اند. پس کوچک‌ترین واحد زمان وجود دارد و تنها لحظه‌های میانی‌ای که می‌توانند مبدأ یا مقصد یک گذار باشند، لحظه‌های صحیح‌اند. بنابراین کافی است هر یال $s\xrightarrow{d}t$ با $d>1$ پیش از اجرای الگوریتم به زنجیرهٔ زیر بازنویسی شود:
\[
s\xrightarrow{1}v_1\xrightarrow{1}v_2\xrightarrow{1}\cdots
\xrightarrow{1}v_{d-1}\xrightarrow{1}t.
\]
حالت‌های $v_i$ همان پیکربندی‌های میانی‌اند که ویژگی پیوستگی وجودشان را تضمین می‌کند. شناسه‌های این حالت‌ها تازه و یکتا انتخاب می‌شوند تا با حالت‌های اظهارشده در ورودی اشتباه نشوند. پس از این بازنویسی، مثال بالا به‌درستی حل می‌شود: در $B$ حالت میانی متناظر با $q_1$ ساخته شده و رابطه کامل می‌شود.

این بازنویسی یک سود الگوریتمی دیگر هم دارد. پس از آن، تنها برچسب زمانی موجود در مدل، تأخیر یک‌واحدی است و رسیدن ضعیف با مدت $d$ برابر با ترکیب $d$بارهٔ رسیدن ضعیف یک‌واحدی است. در نتیجه مجموعهٔ برچسب‌هایی که پالایش افراز باید بررسی کند متناهی و کوچک می‌ماند و الگوریتم لازم نیست همهٔ مدت‌های قابل دستیابی را جداگانه تولید کند. به بیان دیگر، گذر زمان و اجرای کنش پس از این بازنویسی برای الگوریتم افراز یکسان به نظر می‌رسند و هر دو با یک سازوکار بررسی می‌شوند.

\subsection{هزینه و سقف کنترلی}

هزینهٔ این روش افزودن
\[
\sum_{e\in E_{time}}(d_e-1)
\]
حالت میانی است؛ یعنی اندازهٔ کار مستقیماً به بزرگی ثابت‌های زمانی مدل وابسته است. ابزار برای تعداد حالت‌های افزوده‌شده سقفی قابل تنظیم دارد و در صورت عبور از آن، به‌جای رشد کنترل‌نشدهٔ حافظه، اجرا را با پیام خطا متوقف می‌کند. این راه‌حل برای مدل‌های ارزیابی حاضر کافی است، اما در مدل‌هایی با ثابت‌های زمانی بزرگ می‌تواند عامل اصلی مصرف حافظه باشد. فصل \ref{chap:conclusion} راهی را مطرح می‌کند که همین نتیجه را بدون شکستن یال‌ها به دست می‌آورد.

\subsection{دو حالت اجرایی و تفاوت معنایی آن‌ها}

پیاده‌سازی این تصمیم را پنهان نمی‌کند و آن را به‌صورت گزینهٔ \lr{--time-semantics} در دسترس می‌گذارد:

\begin{itemize}
\item حالت \lr{unit}، که پیش‌فرض است و همهٔ نتایج این پایان‌نامه بر مبنای آن گزارش می‌شوند. در این حالت یال‌های تأخیر به‌شرح بالا شکسته می‌شوند و آنچه محاسبه می‌شود دقیقاً افراز متناظر با تعریف \ref{def:weak-timed-bisimulation} است.
\item حالت \lr{strict}، که یال‌های تأخیر را دست‌نخورده نگه می‌دارد و فقط مدت‌های ثبت‌شده در فایل ورودی را به‌عنوان برچسب زمانی می‌پذیرد.
\end{itemize}

باید تصریح کرد که نتیجهٔ حالت \lr{strict} یک شبیه‌سازی دوسویهٔ ضعیف زمان‌دار به معنای تعریف \ref{def:weak-timed-bisimulation} نیست و به‌کاربردن این نام برای آن گمراه‌کننده است. آن حالت در عمل رابطهٔ گذار ورودی را همان‌گونه که ذخیره شده است می‌پذیرد و ویژگی پیوستگی تعریف \ref{def:tts-properties} را به کار نمی‌گیرد؛ یعنی روی ساختاری کار می‌کند که ویژگی‌های استاندارد یک سامانهٔ انتقال زمان‌دار را ندارد. دلیل نگه‌داشتن آن صرفاً تشخیصی است: با آن می‌توان نشان داد که کدام نتیجه‌ها واقعاً به شکستن یال‌ها وابسته‌اند و بنابراین تصمیم معنایی پروژه در کد پنهان نمانده و قابل بازبینی است. بخش \ref{sec:semantic-sensitivity} این تفاوت را روی نمونه‌های واقعی اندازه می‌گیرد. هدف و قرارداد معنایی پروژه در همه حال حالت \lr{unit} است.

خانوادهٔ \lr{mood} در فصل \ref{chap:results} همین وضعیت را در یک مدل واقعی ربکا نشان می‌دهد: در یکی از دو مدل، یک انتظار پنج‌واحدی با یک یال ثبت شده و در دیگری همان انتظار به‌صورت دو واحد، یک گام داخلی و سه واحد ظاهر شده است.


\section{ساخت رابطهٔ گذار ضعیف}

پس از پالایش زمانی، برای هر حالت ابتدا بستار $\tau$ با یک پیمایش عمقی محاسبه می‌شود. سپس برای هر برچسب قابل مشاهده و برچسب تأخیر واحد، مجموعهٔ مقصدهای $\tau^*\ell\tau^*$ ساخته می‌شود. بستار $\tau$ نیز به‌عنوان حرکت تهی در پالایش شرکت می‌کند؛ این همان حالت مدت صفر در تعریف \ref{def:weak-transition} است. حذف این حالت می‌تواند شاخهٔ داخلی منتهی به بن‌بست را از دید الگوریتم پنهان کند.

در پیاده‌سازی، حالت‌ها یک‌بار شماره‌گذاری و مجموعهٔ مقصدها با \lr{BitSet} نگهداری می‌شوند. بنابراین، محاسبهٔ اجتماع مقصدها و مقایسهٔ امضاها از عملیات بیتی استفاده می‌کند. ترتیب حالت‌ها از ورودی و ترتیب برچسب‌ها به‌صورت مرتب‌شده حفظ می‌شود تا اجرای دوباره روی ورودی یکسان مصنوعات متنی یکسانی تولید کند \cite{AfraWeakTimedReduction2026}.

\section{ایدهٔ مبنا: پالایش افراز برای شبیه‌سازی دوسویهٔ انشعابی}
\label{sec:branching-algorithm}

الگوریتم کاهشی که در این پروژه پیاده‌سازی شده است، الگوریتم تازه‌ای از پایه نیست. ساختار آن از الگوریتم شناخته‌شدهٔ گروته و واندراخر برای محاسبهٔ شبیه‌سازی دوسویهٔ انشعابی روی سامانه‌های انتقال برچسب‌دار گرفته شده و برای حالت زمان‌دار تطبیق داده شده است \cite{GrooteVaandrager1990}. پیش از شرح الگوریتم پروژه، لازم است ایدهٔ آن الگوریتم با دقت کافی بیان شود تا روشن شود چه چیزی وام گرفته شده و چه چیزی تغییر کرده است.

\subsection{مسئلهٔ درشت‌ترین افراز}

نقطهٔ شروع این خانواده از الگوریتم‌ها مسئلهٔ \emph{درشت‌ترین افراز} است که پیج و تارجان صورت‌بندی کردند \cite{PaigeTarjan1987}. به‌جای آنکه همهٔ رابطه‌های ممکن میان حالت‌ها جست‌وجو شوند، از یک افراز درشت آغاز می‌شود و تا رسیدن به پایداری شکسته می‌شود. اصطلاح کلیدی، \emph{پایداری} است: افراز $P$ نسبت به بلوک $B'$ و برچسب $a$ پایدار است، هرگاه برای هر بلوک $B\in P$، یا همهٔ اعضای $B$ بتوانند با کنش $a$ به $B'$ برسند، یا هیچ‌کدام نتوانند. اگر بلوکی این شرط را نقض کند، آن بلوک به دو بخشِ «می‌رسند» و «نمی‌رسند» شکسته می‌شود و زوج $(B',a)$ را \emph{شکننده} می‌نامیم.

سه نکته این روش را کارآمد می‌کند. نخست، هر شکستن، افراز را ریزتر می‌کند و چون تعداد حالت‌ها متناهی است، فرایند حتماً پایان می‌یابد. دوم، بلوک‌های شکسته‌شده هرگز دوباره ادغام نمی‌شوند، پس افراز نهایی یکتاست. سوم، افراز نهایی \emph{درشت‌ترین} افراز پایدار است؛ یعنی هیچ دو حالتی را که واقعاً تمایزپذیرند در یک رده نگه نمی‌دارد و هیچ دو حالت تمایزناپذیری را نیز بیهوده جدا نمی‌کند. برای شبیه‌سازی دوسویهٔ قوی، همین افراز دقیقاً رده‌های هم‌ارزی مطلوب است \cite{PaigeTarjan1987,BaierKatoen2008}.

\subsection{آنچه گروته و واندراخر افزودند}

اگر مدل کنش داخلی داشته باشد، شرط بالا بیش از اندازه سخت‌گیر است. یک گام $\tau$ نباید صرفاً به‌دلیل وجودش باعث تمایز دو حالت شود. گروته و واندراخر این مشکل را با مفهوم گذار $\tau$ \emph{بی‌اثر} حل کردند: گذار $s\xrightarrow{\tau}t$ در افراز جاری بی‌اثر است، هرگاه $s$ و $t$ در یک بلوک باشند. چنین گامی ردهٔ رفتاری را تغییر نداده است و الگوریتم آن را نادیده می‌گیرد؛ این همان بند اول تعریف \ref{def:branching-bisimulation} است \cite{GrooteVaandrager1990,GlabbeekWeijland1996}.

با این مفهوم، شرط پایداری بازنویسی می‌شود: بلوک $B$ نسبت به زوج $(B',a)$ پایدار است، هرگاه یا همهٔ اعضای $B$ بتوانند با دنباله‌ای از گذارهای $\tau$ بی‌اثر—یعنی بدون خروج از $B$—و سپس یک گذار $a$ به $B'$ برسند، یا هیچ‌کدام نتوانند. مسیرهای داخلیِ مجاز در پاسخ، همان مسیرهایی‌اند که ردهٔ رفتاری را ترک نمی‌کنند. هر بار که بلوکی شکسته می‌شود، ممکن است گذارهایی که تا پیش از آن بی‌اثر بودند دیگر بی‌اثر نباشند؛ ازاین‌رو الگوریتم باید بلوک‌های متأثر را دوباره بررسی کند. این تکرار تا جایی ادامه می‌یابد که هیچ زوج شکننده‌ای باقی نماند. الگوریتم حاصل پیچیدگی زمانی $O(n(n+m))$ و پیچیدگی فضایی $O(n+m)$ دارد که $n$ تعداد حالت‌ها و $m$ تعداد گذارهاست \cite{GrooteVaandrager1990}؛ نسخه‌های بعدی این مرز را به $O(m\log n)$ رسانده‌اند \cite{GrooteEtAl2017}.

\subsection{آنچه در این پروژه تغییر کرده است}

الگوریتم این پروژه سه عنصر بالا را نگه می‌دارد: شروع از یک افراز درشت، شکستن بلوک‌ها بر اساس بلوک‌های قابل دستیابی، و ادامه تا رسیدن به پایداری. سه تفاوت اصلی نیز دارد:

\begin{enumerate}
\item \textbf{هم‌ارزی هدف.} معیار این پروژه شبیه‌سازی دوسویهٔ \emph{ضعیف} زمان‌دار است، نه انشعابی. بنابراین شرط عضویت حالت‌های میانی مسیر داخلی در همان رده—یعنی شرط $(p,q_1)\in R$ در تعریف \ref{def:branching-bisimulation}—برداشته می‌شود و به‌جای مسیرهای $\tau$ بی‌اثر، بستار کامل $\tau$ به کار می‌رود. این تصمیم از تعریف \ref{def:weak-timed-bisimulation} می‌آید که مبنای پذیرش پروژه است، نه از انتخاب الگوریتمی.
\item \textbf{حضور زمان در مجموعهٔ برچسب‌ها.} در مدل زمان‌دار، برچسب یک گذار می‌تواند مدت نیز باشد. پس از بازنویسی بخش \ref{sec:unit-splitting}، تنها برچسب زمانی موجود تأخیر یک‌واحدی است؛ در نتیجه زمان نیز مانند یک کنش معمولی در همان سازوکار پایداری شرکت می‌کند و لازم نیست برای هر مدت ممکن، یک برچسب جداگانه ساخته شود.
\item \textbf{شکستن مبتنی بر امضا.} به‌جای انتخاب یک زوج شکننده در هر دور، این پیاده‌سازی برای هر حالت یک امضا از مجموعهٔ بلوک‌های قابل دستیابی می‌سازد و همهٔ بلوک‌ها را در یک گذر بر اساس امضا می‌شکند. این صورت ساده‌تر، همان افراز پایدار را می‌دهد اما به مرز پیچیدگی بهینهٔ منابع بالا نمی‌رسد؛ برای اندازهٔ مدل‌های این پروژه کافی است و بخش \ref{chap:conclusion} بهبود آن را در کارهای آینده مطرح می‌کند.
\end{enumerate}


\section{پالایش افراز}

الگوریتم زیر همان طرح بخش \ref{sec:branching-algorithm} است که با سه تغییر یادشده برای شبیه‌سازی دوسویهٔ ضعیف زمان‌دار بازنویسی شده است. الگوریتم با افراز یک‌بلوکی $P_0=\{S\}$ آغاز می‌شود. برای افراز جاری $P$ و برچسب $\ell$، مجموعهٔ بلوک‌های قابل دستیابی از $s$ برابر است با
\[
R_P(s,\ell)=\{[t]_P\mid s\overset{\ell}{\Rightarrow}_W t\}.
\]
امضای هر حالت از بلوک فعلی آن، مجموعهٔ بلوک‌های قابل دستیابی برای تمام برچسب‌ها، و بلوک‌های موجود در بستار $\tau$ ساخته می‌شود:
\[
\operatorname{Sig}_P(s)=
\Bigl([s]_P,\bigl(R_P(s,\ell)\bigr)_{\ell\in L},
\{[t]_P\mid t\in\tau^*(s)\}\Bigr).
\]
حالت‌های یک بلوک که امضای متفاوت دارند از یکدیگر جدا می‌شوند. وجود $[s]_P$ در امضا تضمین می‌کند دورهای بعدی فقط بلوک‌های موجود را بشکنند و دو بلوک جداشده دوباره ادغام نشوند. فرایند زمانی پایان می‌یابد که هیچ امضای تازه‌ای باعث شکستن بلوک نشود. افراز حاصل پایدار است و رده‌های شبیه‌سازی دوسویهٔ ضعیف زمان‌دار را برای معنای ورودی محاسبه می‌کند.

الگوریتم‌های \ref{alg:tau-closure} تا \ref{alg:partition-refinement} شبه‌کد اصلی این فرایند را در چهار مرحله نشان می‌دهند. جداسازی این مراحل روشن می‌کند که رابطهٔ ضعیف چگونه از گذارهای مستقیم ساخته می‌شود و افراز در هر دور دقیقاً بر چه اساسی شکسته می‌شود. در این الگوریتم‌ها، $C_\tau(s)$ بستار کنش داخلی از حالت $s$، $L$ مجموعهٔ برچسب‌های غیر از $\tau$ و $W$ رابطهٔ گذار ضعیف است. در حالت پیش‌فرض پیاده‌سازی، پالایش واحدی پیش از اجرای این مراحل انجام می‌شود و تنها برچسب زمانی موجود در $L$ برابر یک واحد است؛ در خوانش سخت‌گیرانه، $L$ مدت‌های ثبت‌شده روی یال‌های ورودی را نیز در بر می‌گیرد.

الگوریتم \ref{alg:tau-closure} بستار داخلی یک حالت را محاسبه می‌کند. این بستار شامل همهٔ حالت‌هایی است که با صفر یا چند گذار داخلی دسترس‌پذیرند. افزودن خود $s$ به مجموعهٔ آغازین، حرکت داخلی تهی را نیز پوشش می‌دهد. مجموعهٔ \lr{visited} مانع بررسی دوبارهٔ حالت‌ها می‌شود و پایان پیمایش را در حضور چرخه‌های $\tau$ تضمین می‌کند.

\begin{algorithm}[htbp]
\onehalfspacing
\caption{\lr{Tau Closure}}
\label{alg:tau-closure}
\begin{latin}
\begin{algorithmic}[1]
\REQUIRE A transition system $T=(S,L\cup\{\tau\},\rightarrow)$ and a state $s\in S$
\ENSURE All states reachable from $s$ through zero or more $\tau$ transitions
\STATE $\mathit{visited}\gets\{s\}$
\STATE $\mathit{stack}\gets[s]$
\WHILE{$\mathit{stack}$ is not empty}
    \STATE $x\gets\operatorname{pop}(\mathit{stack})$
    \FOR{each transition $(x,\tau,y)\in\rightarrow$}
        \IF{$y\notin\mathit{visited}$}
            \STATE $\mathit{visited}\gets\mathit{visited}\cup\{y\}$
            \STATE $\operatorname{push}(\mathit{stack},y)$
        \ENDIF
    \ENDFOR
\ENDWHILE
\RETURN $\mathit{visited}$
\end{algorithmic}
\end{latin}
\end{algorithm}

الگوریتم \ref{alg:weak-successors} یک برچسب غیر از $\tau$ را میان دو بستار داخلی قرار می‌دهد. اگر $\ell$ یک کنش قابل مشاهده باشد، خروجی همان $\tau^*\ell\tau^*$ است؛ اگر $\ell$ یک برچسب زمانی باشد، خروجی مقصدهای گذار ضعیف با همان مدت را می‌دهد. در حالت \lr{unit}، تکرار این رابطهٔ یک‌واحدی مدت‌های بزرگ‌تر را می‌سازد.

\begin{algorithm}[htbp]
\onehalfspacing
\caption{\lr{Weak Successors for One Label}}
\label{alg:weak-successors}
\begin{latin}
\begin{algorithmic}[1]
\REQUIRE A transition system $T$, tau closures $C_\tau$, a state $s$, and a label $\ell\neq\tau$
\ENSURE The targets of the weak step $s\overset{\ell}{\Rightarrow}_W t$
\STATE $\mathit{targets}\gets\varnothing$
\FOR{each state $x\in C_\tau(s)$}
    \FOR{each transition $(x,\ell,y)\in\rightarrow$}
        \FOR{each state $t\in C_\tau(y)$}
            \STATE $\mathit{targets}\gets\mathit{targets}\cup\{t\}$
        \ENDFOR
    \ENDFOR
\ENDFOR
\RETURN $\mathit{targets}$
\end{algorithmic}
\end{latin}
\end{algorithm}

الگوریتم \ref{alg:weak-relation} دو محاسبهٔ قبلی را برای همهٔ حالت‌ها و برچسب‌ها انجام می‌دهد. خروجی آن علاوه بر $W$، بستارهای داخلی را نیز نگه می‌دارد؛ زیرا $C_\tau(s)$ حالت مدت صفر تعریف \ref{def:weak-transition} است و در پالایش افراز باید مانند دیگر حرکت‌های ضعیف بررسی شود.

\begin{algorithm}[htbp]
\onehalfspacing
\caption{\lr{Weak Transition Relation}}
\label{alg:weak-relation}
\begin{latin}
\begin{algorithmic}[1]
\REQUIRE A timed transition system $T=(S,\mathit{Act}\cup\{\tau\}\cup D,\rightarrow)$
\ENSURE Tau closures $C_\tau$, non-tau labels $L$, and weak relation $W$
\FOR{each state $s\in S$}
    \STATE $C_\tau(s)\gets\operatorname{TauClosure}(T,s)$
\ENDFOR
\STATE $L\gets\{\ell\mid (x,\ell,y)\in\rightarrow\ \land\ \ell\neq\tau\}$
\STATE $W\gets\varnothing$
\FOR{each state $s\in S$}
    \FOR{each label $\ell\in L$}
        \STATE $A\gets\operatorname{WeakSuccessors}(T,C_\tau,s,\ell)$
        \FOR{each state $t\in A$}
            \STATE $W\gets W\cup\{(s,\ell,t)\}$
        \ENDFOR
    \ENDFOR
\ENDFOR
\RETURN $(C_\tau,L,W)$
\end{algorithmic}
\end{latin}
\end{algorithm}

الگوریتم \ref{alg:partition-refinement} با قرار دادن همهٔ حالت‌ها در یک بلوک آغاز می‌شود. در هر دور، برای هر حالت مجموعهٔ بلوک‌های قابل دسترس با هر برچسب و مجموعهٔ بلوک‌های موجود در بستار داخلی محاسبه می‌شود. دو حالت تنها زمانی در یک بلوک می‌مانند که بلوک فعلی و همهٔ این مجموعه‌ها برای آن‌ها یکسان باشند. مقایسهٔ آرایهٔ $\mathit{nextBlock}$ با $\mathit{block}$ شرط توقف را به‌صورت صریح نشان می‌دهد.

\begin{algorithm}[htbp]
\onehalfspacing
\caption{\lr{Stable Partition Refinement}}
\label{alg:partition-refinement}
\begin{latin}
\begin{algorithmic}[1]
\REQUIRE States $S$, tau closures $C_\tau$, labels $L$, and weak relation $W$
\ENSURE The coarsest stable partition $P$
\FOR{each state $s\in S$}
    \STATE $\mathit{block}[s]\gets 0$
\ENDFOR
\REPEAT
    \STATE $\mathit{signatureIds}\gets$ an empty ordered map
    \FOR{each state $s\in S$}
        \FOR{each label $\ell\in L$}
            \STATE $R_\ell(s)\gets\{\mathit{block}[t]\mid(s,\ell,t)\in W\}$
        \ENDFOR
        \STATE $R_{\tau^*}(s)\gets\{\mathit{block}[t]\mid t\in C_\tau(s)\}$
        \STATE $\sigma(s)\gets(\mathit{block}[s],(R_\ell(s))_{\ell\in L},R_{\tau^*}(s))$
        \IF{$\sigma(s)\notin\mathit{signatureIds}$}
            \STATE Assign the next unused block id to $\mathit{signatureIds}[\sigma(s)]$.
        \ENDIF
        \STATE $\mathit{nextBlock}[s]\gets\mathit{signatureIds}[\sigma(s)]$
    \ENDFOR
    \STATE $\mathit{changed}\gets(\mathit{nextBlock}\neq\mathit{block})$
    \STATE $\mathit{block}\gets\mathit{nextBlock}$
\UNTIL{$\mathit{changed}=\mathit{false}$}
\STATE $P\gets\bigl\{\{s\in S\mid\mathit{block}[s]=i\}\mid i\text{ is an assigned id}\bigr\}$
\RETURN $P$
\end{algorithmic}
\end{latin}
\end{algorithm}

پیاده‌سازی مبتنی بر امضا در هر دور، برای هر حالت و هر برچسب مجموعهٔ بلوک‌های مقصد را بررسی می‌کند و حداکثر به تعداد حالت‌ها دور دارد. الگوریتم‌های پالایش پیشرفته‌تری با کران مجانبی بهتر وجود دارند، اما روش حاضر با شبه‌کد پروژه تطابق مستقیم دارد و برای بزرگ‌ترین ورودی موجود، که ۴۲ حالت پیش از پالایش زمانی دارد، کافی بوده است. فصل \ref{chap:results} از این اندازه فراتر ادعای مقیاس‌پذیری نمی‌کند.

\section{ساخت مدل خارج‌قسمتی و فشرده‌سازی زمان}

در ساخت خارج‌قسمت، هر ردهٔ هم‌ارزی به یک حالت تبدیل می‌شود و اگر یالی با برچسب $\ell$ از عضوی از ردهٔ $B$ به عضوی از ردهٔ $B'$ وجود داشته باشد، یال $B\xrightarrow{\ell}B'$ ساخته می‌شود. خودحلقهٔ $\tau$ درون یک رده حذف می‌شود؛ زیرا چیزی به بستار داخلی آن رده اضافه نمی‌کند.

خارج‌قسمت مستقیم سامانهٔ پالایش‌شده فقط یال‌های یک‌واحدی دارد. در نتیجه، یک انتظار ده‌واحدی می‌تواند چندین ردهٔ میانی ایجاد کند و حتی مدل ظاهراً کاهش‌یافته را از ورودی بزرگ‌تر سازد. برای جلوگیری از این مسئله، زنجیرهٔ بیشینه‌ای از رده‌ها که تنها یک یال زمانی ورودی و یک یال زمانی خروجی دارند، حالت آغازین نیستند و هیچ انتخاب یا کنش قابل مشاهده‌ای ندارند، به یک یال با مدت مجموع تبدیل می‌شود. شرط ساختاری سخت‌گیرانهٔ این مرحله مانع حذف نقطه‌ای می‌شود که محل انتخاب رفتاری است.

به همین دلیل، تعداد رده‌های افراز لزوماً با تعداد حالت‌های خروجی یکسان نیست. هر رده یا به یک حالت خارج‌قسمتی تبدیل می‌شود یا موقعیت آن روی یک یال تأخیر فشرده‌شده ثبت می‌شود. فایل \lr{partition.json} این نگاشت را نگه می‌دارد. همچنین، اطلاعات لازم برای بازکردن دوبارهٔ یال‌های چندواحدی ذخیره می‌شود تا وارسی‌گر بتواند دقیقاً مدل تحویلی را با سامانهٔ پالایش‌شده مقایسه کند.

\section{وارسی مستقل خروجی}

وارسی‌گر رابطهٔ میان حالت‌های سامانهٔ پالایش‌شده و مدل خارج‌قسمتی را مستقیماً در برابر تعریف شبیه‌سازی دوسویهٔ ضعیف زمان‌دار می‌سنجد. برای هر زوج مرتبط، همهٔ حرکت‌های تهی و همهٔ گذارهای ضعیف یک طرف باید با گذار هم‌برچسب طرف دیگر پاسخ داده شوند و مقصدها نیز مرتبط بمانند؛ این بررسی در هر دو جهت انجام می‌شود. حالت‌های آغازین نیز باید در رابطه باشند.

وارسی‌گر از بستارهای محاسبه‌شده در کاهش‌دهنده استفاده نمی‌کند. برای نمونه، بستارها را با مجموعه‌های معمولی از نو می‌سازد؛ درحالی‌که کاهش‌دهنده مقصدها را در ساختار \lr{BitSet} نگه می‌دارد. استقلال کامل منطقی ادعا نمی‌شود، زیرا هر دو مؤلفه تعریف معنایی مشترکی دارند. بااین‌حال، جدایی کد باعث می‌شود خطای جزئی در شاخص‌گذاری یا عملیات بیتی کاهش‌دهنده، خودبه‌خود در وارسی‌گر تکرار نشود.

برای نشان‌دادن توان ردکردن، پنج تغییر مخرب روی مدل خارج‌قسمتی اعمال شده است: حذف یال، تغییر برچسب، تغییر مدت، تغییر مقصد و افزودن یال. انتظار می‌رود وارسی‌گر همهٔ این تغییرها را رد کند.

\section{معماری نرم‌افزار و رابط اجرا}

پیاده‌سازی با جاوای ۱۷ و \lr{Maven} ساخته شده است. نسخهٔ جاوا و نسخه‌های افزونه‌های ساخت با مخازن رسمی کامپایلر و وارسی‌گر مدل ربکا هم‌راستا شده‌اند \cite{RebecaCompilerGitHub2026,RebecaModelCheckerGitHub2026}. وابستگی زمان اجرا خارج از \lr{JDK} وجود ندارد؛ خواندن فایل با \lr{StAX} و تولید \lr{JSON} با نویسنده‌ای کوچک و قطعی انجام می‌شود. این انتخاب احتمال تعارض کتاب‌خانه‌ها را در صورت ادغام آینده با محیط مبتنی بر \lr{Eclipse} افرا کاهش می‌دهد.

بسته‌های اصلی نرم‌افزار مسئولیت‌های زیر را دارند:

\begin{itemize}
\item \lr{afra} و \lr{source}: خواندن فایل و تبدیل منبع داده به نمایش خام؛
\item \lr{tts}: نمایش تغییرناپذیر و اعتبارسنجی‌شدهٔ سامانهٔ انتقال؛
\item \lr{weak}: پالایش زمانی، ساخت گذارهای ضعیف و پالایش افراز؛
\item \lr{quotient}: ساخت خارج‌قسمت و فشرده‌سازی زنجیره‌های زمانی؛
\item \lr{verify}: بررسی مستقل رابطهٔ خروجی؛
\item \lr{report}: تولید مصنوعات و سنجه‌ها؛
\item \lr{app}: هماهنگ‌کردن مراحل و رابط خط فرمان.
\end{itemize}

\gls{CommandLineInterface} چهار دستور دارد: \lr{inspect} برای شناخت محتوای ورودی، \lr{reduce} برای کاهش یک سامانه، \lr{equivalent} برای مقایسهٔ حالت آغازین دو سامانه، و \lr{visualize} برای تولید نمایش \lr{Graphviz}. منطق اصلی در \lr{ReductionService} قرار دارد و خط فرمان تنها گزینه‌ها و کد خروج را مدیریت می‌کند؛ بنابراین، همان جریان بدون ایجاد فرایند جداگانه از کد جاوا نیز قابل فراخوانی است. کد خروج صفر موفقیت، یک نتیجهٔ معنایی منفی، دو ورودی نامعتبر و سه نقض یک ناوردای داخلی را نشان می‌دهد.

\section{مصنوعات خروجی}

جدول \ref{tab:output-artifacts} نقش پنج فایل تولیدشده برای هر کاهش را نشان می‌دهد. هیچ‌یک از این فایل‌ها به‌تنهایی جای همهٔ بقیه را نمی‌گیرد؛ برای نمونه، نمایش سازگار با افرا برای مشاهده مناسب است، ولی عضویت رده‌ها و منشأ کامل یال‌ها را نگه نمی‌دارد.

\begin{table}[htbp]
\caption{مصنوعات تولیدشده برای هر اجرای کاهش}
\label{tab:output-artifacts}
\centering
\onehalfspacing
\begin{tabularx}{0.94\textwidth}{|l|X|}
\hline
فایل & محتوا و کاربرد \\
\hline
\lr{reduced.json} & نمایش اصلی و بدون اتلاف مدل کاهش‌یافته، شامل برچسب‌ها، اعضای رده‌ها و مجموعهٔ مشاهده‌ها. \\
\hline
\lr{partition.json} & عضویت هر حالت ورودی در رده و موقعیت آن در مدل کاهش‌یافته. \\
\hline
\lr{reduced.statespace} & نمای سازگار با ابزار نمایش فضای حالت افرا؛ این نما به‌دلیل محدودیت طرح داده اتلاف‌دار است. \\
\hline
\lr{reduced.dot} & گراف \lr{DOT} برای ترسیم ساختار مدل خارج‌قسمتی. \\
\hline
\lr{metrics.json} & اندازه‌ها، نسبت‌های کاهش، زمان و حافظه، نتیجهٔ وارسی و اطلاعات ردیابی اجرا. \\
\hline
\end{tabularx}
\end{table}

قطعی‌بودن ترتیب حالت‌ها، رده‌ها و یال‌ها موجب می‌شود چهار مصنوع ساختاری در دو اجرای یکسان، بایت‌به‌بایت برابر باشند. زمان اجرا و بیشینهٔ حافظه طبیعتاً در \lr{metrics.json} متغیرند. هش \lr{SHA-256} ورودی و شناسهٔ کامیت امکان اتصال هر نتیجه به فایل و نسخهٔ مولد آن را فراهم می‌کند.

\section{طرح اعتبارسنجی و تکرار اجرا}

اعتبارسنجی در سه سطح انجام شده است. در سطح معنایی، نمونه‌های کوچک با افراز مورد انتظار، بستار $\tau$، مشاهده‌ها، تأخیرها، بن‌بست و ورودی نامعتبر را می‌سنجند. در سطح مدل خارج‌قسمتی، کامل‌بودن افراز، معتبر‌بودن مقصد یال‌ها، رابطهٔ حالت آغازین و رد جهش‌های عمدی بررسی می‌شود. در سطح انتهابه‌انتها، فایل \lr{.statespace} به برنامهٔ بسته‌بندی‌شده داده می‌شود و کد خروج، شمار حالت‌ها و محتوای همهٔ مصنوعات کنترل می‌شوند. جزئیات عددی این آزمون‌ها در فصل \ref{chap:results} آمده است.

برای تکرار ساخت و آزمون نسخهٔ ثبت‌شدهٔ پروژه، فرمان‌های اصلی به‌شکل زیرند:

\begin{latin}
\begin{verbatim}
mvn clean package
python3 tests/e2e/run_e2e.py
python3 evaluation/run_evaluation.py
\end{verbatim}
\end{latin}

فرمان نخست نرم‌افزار و آزمون‌های واحد را می‌سازد، فرمان دوم رابط بسته‌بندی‌شده را مانند یک کاربر خارجی می‌آزماید و فرمان سوم فایل جمع‌بندی و سنجه‌های خام ارزیابی را از نو تولید می‌کند. ورودی‌ها، اسکریپت‌ها و نتایج ثبت‌شده در مخزن پروژه موجودند \cite{AfraWeakTimedReduction2026}.

\section{جمع‌بندی}

روش این پروژه یک زنجیرهٔ قابل ردیابی از خروجی \lr{TTS} افرا تا مدل خارج‌قسمتی است. مجموعهٔ مشاهده‌ها صریح است، تفاوت دو تفسیر تأخیر به گزینه و آزمون تبدیل شده، پالایش افراز از قرارداد ورودی جدا مانده و مدل تحویلی با پیاده‌سازی دیگری از شروط انتقال بررسی می‌شود. فصل بعد نشان می‌دهد این تصمیم‌ها روی نمونه‌های ارزیابی چه خروجی‌ای تولید کرده‌اند و شواهد تا چه اندازه از درستی و کارایی روش پشتیبانی می‌کنند.
